\chapter*{Introduction to Part~II}
\addcontentsline{toc}{chapter}{Introduction to Part~II}
\label{ch:part2-intro}

In this Part we turn from the physics of the inflationary epoch to the present-day dark sector, investigating the nature of dark matter and dark energy through a combination of cosmological and terrestrial observations. Dark matter and dark energy together account for roughly 95\% of the energy budget of the Universe, yet their fundamental nature remains unknown. The four chapters that follow attack this problem from complementary directions, unified by a common methodological theme of interplay between theoretical modeling and precision observational data in constraining the properties of the dark sector.

We begin with two studies of particle dark matter production and its observable consequences. Chapter~\ref{ch:minimal-dm-freezein} explores dark matter production via freeze-in through an ultralight dark photon portal, focusing on a regime that has received little attention in the literature in which the reheating temperature lies below the dark matter mass. In this regime the production is Boltzmann-suppressed, requiring larger portal couplings to reproduce the observed relic abundance and thereby lifting the freeze-in benchmark into the reach of current direct detection experiments such as SENSEI, DAMIC-M, and XENON. Chapter~\ref{ch:cmb-leptophilic-dm} investigates sub-MeV leptophilic dark matter that couples to Standard Model leptons and finds that the dominant observable signature is not the tree-level scattering but rather a loop-induced decay into two photons. Although the resulting lifetimes are cosmologically long, the associated energy injection leaves imprints on CMB anisotropies that we constrain using the latest \textit{Planck} NPIPE data, obtaining bounds that surpass existing terrestrial limits by up to two orders of magnitude in parts of the parameter space.

The second half of this Part broadens the scope beyond particle dark matter. Chapter~\ref{ch:pbh-memory-burden} examines whether the recently proposed memory-burden effect --- a mechanism that could suppress Hawking evaporation by stabilizing black holes through the information they carry --- can open a new mass window for primordial black holes as dark matter in the range $10^4$--$10^{10}\,$g. By confronting this scenario with CMB, BBN, and gamma-ray constraints we show that such a window is viable only if the transition to the memory-burdened phase is nearly instantaneous, requiring the black hole to lose essentially no mass before stabilization occurs. Finally, Chapter~\ref{ch:desi-dark-energy} addresses the dark energy side of the dark sector, critically examining the recent DESI DR2 evidence for time-varying dark energy. Using theory-informed priors derived from well-motivated quintessence models and a normalizing-flow framework for efficient posterior reweighting, we demonstrate that the apparent ${\sim}3\sigma$ preference for dynamical dark energy is substantially reduced once physically motivated priors replace the standard uniform ones, highlighting the sensitivity of marginal detections to prior assumptions.
